\chapter{Communicating Risk Across the Organisation}
\label{chap:Communicating Risk Across the Organisation}

%---------------------------- SECTION ----------------------------
\section{The Gap Between Knowing and Understanding}
\paragraph{There is a frustrating paradox at the heart of risk management in many organisations. The risks are identified. The assessments are conducted. The registers are maintained. The reports are written. And yet, somewhere between the compliance function producing that work and the people who need to act on it actually changing their behaviour, something is lost. The message does not land. The report is read but not absorbed. The risk rating is noted but its implications are not understood. The board approves the risk appetite statement without genuinely engaging with what it means in practice.}
\paragraph{This is the communication problem — and it is more prevalent, more consequential, and more underestimated than almost any other challenge in risk management.}
\paragraph{Risk communication is not simply the act of distributing information about risk. It is the process of ensuring that the right people develop a genuine understanding of the risks that are relevant to them, that they appreciate the significance of those risks in a way that actually influences their decisions and behaviour, and that the organisation develops and maintains a shared, accurate picture of its risk landscape. That is a considerably more demanding task than writing a report or updating a register — and it requires a different set of skills, a different kind of thinking, and a considerably greater investment of attention than most organisations give it.}
\paragraph{For compliance and legal professionals, risk communication is a core professional competency. The quality of the risk assessments produced means nothing if the insights they contain cannot be conveyed effectively to the people who need them. The ability to communicate risk clearly, honestly, and persuasively — to a board, to a senior management team, to operational colleagues, or to a regulator — is one of the most practically valuable skills in the compliance professional's toolkit.}
\paragraph{This chapter explores why risk communication is difficult, what makes it effective, and how to tailor it to the very different audiences and purposes that compliance professionals typically encounter.}

%---------------------------- SECTION ----------------------------
\section{Why Risk Communication Is Hard}

\paragraph{Before exploring how to communicate risk well, it is worth understanding why it is genuinely difficult — because the challenges are not merely technical. They are deeply rooted in human psychology, organisational dynamics, and the inherent complexity of risk itself.}

\subsection{Risk Is Abstract Until It Is Not}
\paragraph{One of the fundamental challenges of risk communication is that risk, by definition, concerns things that have not yet happened. Human beings are significantly better at responding to concrete, visible, immediate threats than to abstract, statistical, future ones. A risk described as \textit{"a 15\% probability of a significant regulatory breach within the next twelve months"} is real and meaningful to the person who constructed that estimate — but it is likely to feel abstract and remote to a board member who has not personally experienced a regulatory investigation and for whom the phrase \textit{"significant regulatory breach"} conjures only a vague and distant threat.}
\paragraph{This is not a failure of intelligence or attention on the part of the board member. It is a fundamental feature of human cognition — the same feature that makes people underinvest in pension savings, underestimate the risk of road accidents, and fail to take adequate precautions against events they have never personally experienced. Effective risk communication needs to work with this cognitive reality rather than against it — finding ways to make abstract risks concrete, immediate, and personally relevant to the audience.}

\subsection{Numbers Without Narrative Are Not Communication}
\paragraph{Risk assessments are full of numbers — likelihood ratings, consequence scores, risk ratings, KRI thresholds, financial loss estimates. These numbers carry genuine analytical content, but presented without context, without narrative, and without explanation of what they mean and why they matter, they communicate very little. A board that is told that fourteen risks are currently rated red, compared to eleven last quarter, knows that something has changed — but has no basis for understanding whether this represents a serious deterioration, a normal fluctuation, or simply a change in the rating methodology.}
\paragraph{Numbers tell you what. Narrative tells you why, what it means, and what should be done about it. Both are necessary — and in most risk communication, the narrative does far more of the real communicative work.}

\subsection{The Messenger and the Message}
\paragraph{Risk communication is never purely about the content of the message. It is also about the relationship between the messenger and the audience — the degree of trust, credibility, and mutual understanding that exists, and the way in which the message is positioned within that relationship.}
\paragraph{A compliance professional who is perceived as a business partner — someone who understands the commercial context, who provides proportionate and practical advice, and who raises concerns constructively — is likely to find that their risk communications are received with genuine engagement. A compliance professional who is perceived primarily as a constraint, a bureaucrat, or a bearer of bad news is likely to find that even well-constructed risk communications are received with defensiveness or dismissal.}
\paragraph{This does not mean that the compliance function should soften its messages to maintain relationships — that would be a fundamental failure of the oversight role. It means that investing in relationships, building credibility through quality and consistency, and framing risk communications in ways that respect the audience's intelligence and context is not merely good manners — it is a prerequisite for effective communication.}

\subsection{Organisational Dynamics and the Suppression of Bad News}
\paragraph{In many organisations, there are powerful informal pressures against communicating bad news upwards. People who raise concerns are sometimes perceived as negative or obstructive. Teams that surface problems in their own areas may be seen as admitting failure. Senior leaders who have publicly committed to a strategy or a project may be unreceptive to risk communications that challenge its viability.}
\paragraph{These dynamics create a systematic pressure to soften, delay, or suppress risk communications — and they are one of the most significant contributors to the kind of risk accumulation that eventually produces serious organisational failures. The financial crisis, as we noted in Chapter~\ref{chap:A Brief History of Risk Assessment}, involved not just a failure to identify risks but a pervasive failure to communicate them effectively upward through governance structures.}
\paragraph{For compliance professionals, navigating these dynamics is a genuine professional challenge. It requires courage — the willingness to deliver uncomfortable messages clearly and persistently, even when they are unwelcome. It also requires skill — the ability to frame difficult communications in ways that are heard rather than dismissed, that invite engagement rather than defensiveness, and that make the case for action in terms that resonate with the audience's own priorities and concerns.}

%---------------------------- SECTION ----------------------------
\section{Principles of Effective Risk Communication}
\paragraph{Against this background of genuine difficulty, several principles consistently characterise effective risk communication — regardless of the audience, the medium, or the specific risk being communicated.}

\subsection{Know Your Audience}
\paragraph{The most fundamental principle of any communication is to understand who you are communicating with — their level of knowledge, their primary concerns, their decision-making authority, and the context in which they are receiving the communication. Risk communication is no different, and the variation between audiences in a typical organisation is substantial.}
\paragraph{A board member with a finance background will bring different prior knowledge, different instincts, and different questions to a risk report than one with a legal background, a technology background, or an operational background. An operational manager in a front-line business unit needs different information, at a different level of detail, than a senior risk committee member. A regulator reviewing the organisation's risk management arrangements has different expectations and different concerns than an internal audit committee.}
\paragraph{Effective risk communication starts with a genuine effort to understand the audience — what they already know, what they need to know, what decisions they are making, and what form of communication is most likely to engage their genuine attention. This is not about telling people what they want to hear. It is about framing what they need to hear in a way that they can genuinely receive and act upon.}

\subsection{Lead With Significance}
\paragraph{In any risk communication, the most important thing is to be clear about what matters most — to lead with the most significant risks, the most important changes, and the most urgent decisions, rather than burying them in comprehensive coverage of everything.}
\paragraph{This principle runs counter to a natural instinct in risk management — the desire to demonstrate thoroughness by covering everything equally. Comprehensive coverage has its place in detailed working documentation. But in communication to decision-makers, comprehensiveness without prioritisation is a form of noise — it makes it harder, not easier, to identify and engage with the things that actually require attention.}
\paragraph{The discipline of leading with significance — of identifying the two or three things that most need to be understood and acted upon, and making those the centrepiece of the communication — is one of the most valuable skills a compliance professional can develop.}

\subsection{Make the Abstract Concrete}
\paragraph{As noted above, abstract risk descriptions are difficult to engage with. Making risks concrete — using specific scenarios, real examples, analogies, and illustrations of what a risk actually looks like when it materialises — significantly increases the likelihood that the audience will genuinely understand and engage with the risk.}
\paragraph{This does not require the communication to be alarmist or to exaggerate risks for effect. It requires the communicator to think carefully about how to translate statistical or abstract risk descriptions into forms that connect with the audience's experience and imagination.}
\paragraph{\textit{"There is a moderate probability of a significant regulatory breach in our transaction monitoring process"} is a valid risk description. \textit{"Our transaction monitoring system generated 340 alerts last quarter, of which only 12\% were reviewed within the required timeframe. In comparable firms, this pattern has preceded regulatory enforcement action in several recent cases, including a £X million fine imposed on [peer firm] last year"} is a risk communication that gives the board something concrete to engage with.}

\subsection{Distinguish Between Information and Insight}
\paragraph{There is an important difference between providing information and providing insight. Information tells you what has happened or what the current state is. Insight tells you what it means, why it matters, and what should be done about it.}
\paragraph{Many risk reports are rich in information but limited in insight. They contain extensive data — KRI trends, risk ratings, control test results, incident statistics — but they do not clearly articulate what this data means for the organisation's risk profile, what the implications are for management decisions, or what actions are most urgently needed.}
\paragraph{The transition from information to insight is where the compliance professional's judgement and expertise adds the most value. Anyone can collect and present data. The ability to interpret it — to identify the patterns, the anomalies, the concerning trends, and the genuine cause for comfort — and to communicate those interpretations clearly is what distinguishes genuinely valuable risk communication from a data dump.}

\subsection{Be Honest About Uncertainty}
\paragraph{Risk assessments involve judgement under uncertainty — and effective risk communication should be honest about that fact. Presenting risk evaluations as though they were precise measurements, without acknowledging the assumptions, the data limitations, and the degree of uncertainty involved, misleads the audience and undermines the credibility of the communication when reality diverges from the assessment.}
\paragraph{Honest communication about uncertainty does not mean constantly qualifying every statement to the point of uselessness. It means being clear about where the assessment is based on solid evidence, where it relies on judgement, where the range of possible outcomes is wide, and where the organisation is genuinely uncertain. This kind of honesty is not a weakness — it is a mark of intellectual integrity that builds credibility with sophisticated audiences, including regulators.}

\subsection{Close the Loop — Make the Call to Action Clear}
\paragraph{Every risk communication should leave its audience in no doubt about what is expected of them. Does this communication require a decision? If so, what decision, and what are the options? Does it require the allocation of resources? The commissioning of further investigation? The escalation of a concern? A change in policy or process?.}
\paragraph{Risk communication that simply informs — that presents a picture of the risk landscape without making clear what action is expected — is incomplete. The compliance function's role is not merely to inform but to support and sometimes to drive the decisions and actions needed to manage risk effectively. Making the call to action explicit — clearly and specifically — is part of fulfilling that role.}

%---------------------------- SECTION ----------------------------
\section{Communicating with Different Audiences}
\paragraph{With those principles in mind, let us consider how risk communication should be tailored to the specific audiences that compliance professionals most commonly encounter.}

\subsection{Communicating with the Board}
\paragraph{Board-level risk communication is among the most important and most challenging communication tasks a compliance professional faces. The stakes are high — the board's engagement with risk is fundamental to good governance — and the constraints are significant — board time is limited, board members have diverse backgrounds, and the volume of material competing for their attention is substantial.}
\paragraph{Effective board-level risk communication has several specific characteristics:}
\paragraph{\textbf{It is strategic in focus. The board's role is governance and strategic oversight — not operational management. Board-level risk communication should focus on the organisation's most material risks, its overall risk profile relative to appetite, and the strategic implications of significant risk developments. Operational detail belongs in management-level reporting, not board papers.}}
\paragraph{\textbf{It is concise}. A board risk report that runs to fifty pages of dense text and tables will not be read with genuine engagement, however thorough its content. The discipline of producing concise, focused board risk communication — identifying the most important things and presenting them clearly — is both a communication skill and a governance service.}
\paragraph{\textbf{It is forward-looking}. Boards need to understand not just where the organisation's risks stand today but how they are likely to develop — what risks are emerging, what changes in the environment are likely to affect the risk profile, and what decisions the board may need to make in the period ahead. A purely retrospective risk report, however accurate, does not give the board the information it needs to fulfil its governance responsibilities.}
\paragraph{\textbf{It invites genuine engagement}. The best board risk communications do not simply present conclusions — they frame questions, highlight areas of uncertainty, and invite the board to bring its own judgement and experience to bear. A board that is genuinely engaged in risk oversight is a more effective governance body than one that merely ratifies management's conclusions.}
\paragraph{\textbf{It is honest about problems}. The temptation to present board-level risk communications in the most favourable light — to smooth over concerns, emphasise positives, and avoid raising uncomfortable issues — is understandable but counterproductive. A board that is not told about serious risk concerns cannot fulfil its oversight responsibilities, and a compliance function that withholds or softens material information has failed at its most basic level.}

\subsection{Communicating with Senior Management}
\paragraph{Senior management risk communication sits between the strategic focus of board reporting and the operational detail of front-line management — providing an integrated view of the risk profile that supports management decision-making across functions and business units.}
\paragraph{Key characteristics of effective senior management risk communication include:}
\paragraph{\textbf{Cross-functional integration}. Senior management needs to understand how risks interact across different parts of the business — where risks in one area create exposure in another, where control weaknesses in one function affect the risk profile of the whole organisation, and where the combined effect of multiple moderate risks creates a more serious aggregate exposure than any individual risk would suggest.}
\paragraph{\textbf{Action orientation}. Senior management is in a position to make decisions, allocate resources, and drive operational change. Risk communications at this level should be closely connected to the actions required — identifying specifically what needs to happen, who is responsible for making it happen, and by when.}
\paragraph{\textbf{Appropriate challenge}. The second-line compliance function's role at the senior management level includes providing genuine, independent challenge to the first line's assessment of its own risks. This means being willing to express a different view, to push back on risk ratings that appear optimistic, and to raise concerns that the first line may be underweighting. Doing this with appropriate professionalism — constructively and with supporting reasoning — is an important part of the second-line role.}

\subsection{Communicating with Operational Teams}
\paragraph{Communicating risk to operational teams — the people who own and manage risks in the first line — requires a different approach from senior and board-level communication. The focus is less on the strategic picture and more on the specific, practical implications for the way people do their jobs.}
\paragraph{Effective operational risk communication:}
\paragraph{\textbf{Is specific and relevant}. Operational staff are most engaged by risk information that is directly relevant to their own activities and responsibilities. Generic compliance communications that describe risks in abstract terms or at a level of generality that does not connect with daily work are likely to be received with limited attention.}
\paragraph{\textbf{Is practical and action-oriented}. Operational teams need to understand not just what the risks are but what they are expected to do about them — specifically, concretely, and in terms that are compatible with operational realities.}
\paragraph{\textbf{Creates space for two-way communication}. Some of the most valuable risk information in any organisation sits with the people closest to operational processes — the people who know where the system is unreliable, where the procedure is routinely bypassed, and where the control gap has been known for months but never reported upward. Effective operational risk communication creates genuine opportunities for this information to flow upward — through feedback mechanisms, escalation channels, risk workshops, and a culture in which raising concerns is genuinely welcomed.}
\paragraph{\textbf{Is supported by training and awareness}. Risk communication to operational teams is most effective when it is reinforced by training — not compliance training as a box-ticking exercise, but genuine, relevant, engaging learning that helps people understand why the risks matter and what good risk management looks like in their specific context.}

\subsection{Communicating with Regulators}
\paragraph{Regulatory communication is a specific and high-stakes form of risk communication that deserves particular attention. When communicating with regulators — whether in response to supervisory enquiries, in the course of regulatory examinations, or through formal regulatory submissions — the principles of effective risk communication apply with particular force.}
\paragraph{\textbf{Be transparent and accurate}. Regulators expect and are entitled to an honest, accurate account of the organisation's risk profile and risk management arrangements. Presenting a misleading picture — whether through selective emphasis, incomplete disclosure, or technical accuracy that obscures a substantive problem — is both a regulatory and an ethical failure, and carries significant consequences when discovered.}
\paragraph{\textbf{Demonstrate genuine engagement}. Regulators are sophisticated and experienced — they have seen many risk management frameworks and many regulatory submissions, and they are generally well placed to distinguish between organisations that are genuinely engaged in risk management and those that are going through the motions. Communicating in a way that demonstrates genuine understanding of the risks, honest acknowledgement of weaknesses, and credible plans for improvement is far more effective than presenting a polished but superficial picture.}
\paragraph{\textbf{Self-report proactively}. Where the organisation has identified a significant compliance failure or a material control weakness, proactive and timely disclosure to the relevant regulator — before the issue is identified through supervisory activity — is almost always the better course. Regulators consistently treat organisations that self-report and cooperate more favourably than those who conceal or delay. The reputational and regulatory costs of being found to have withheld information from a regulator are typically far greater than the costs of self-reporting.}
\paragraph{\textbf{Maintain consistency}. Regulatory communications need to be consistent with each other and with the organisation's internal risk documentation. Inconsistencies between what is told to regulators and what is recorded internally are a serious governance red flag — and regulators are increasingly sophisticated in their ability to identify them.}

%---------------------------- SECTION ----------------------------
\section{Risk Perception and Human Psychology — Understanding Your Audience's Mental Models}
\paragraph{Effective risk communication requires not just good communication skills but genuine understanding of how human beings perceive and process risk — because the way that people naturally think about risk is often at odds with the formal, analytical frameworks that risk professionals use.}
\paragraph{Several well-documented features of human risk perception are particularly relevant:}

\subsection{The Availability Heuristic}
\paragraph{As discussed in Chapter~\ref{chap:Step 1 — Identifying Hazards and Risk Sources}, people tend to assess the likelihood of an event by how easily examples come to mind — giving disproportionate weight to recent, vivid, or emotionally salient events. In a communication context, this means that risks associated with recent incidents or prominent external events will feel more real and more urgent to your audience than risks that are statistically more significant but less vivid.}
\paragraph{A communicator who understands this can use it constructively — anchoring abstract risk descriptions to concrete examples and analogies that make the risk feel real — without manipulating the audience or misrepresenting the actual risk level.}

\subsection{Loss Aversion}
\paragraph{Behavioural economics has demonstrated that people are significantly more motivated by the prospect of avoiding a loss than by the prospect of achieving an equivalent gain. In risk communication terms, this means that framing risks in terms of what could be lost — regulatory permissions, customer trust, financial resources, personal reputation — is often more motivating than framing them in terms of what could be achieved through better risk management.}
\paragraph{This is not about being alarmist. It is about understanding that risk, by its nature, is about potential loss — and communicating it in terms that connect with the audience's natural tendency to prioritise loss avoidance.}

\subsection{The Optimism Bias Revisited}
\paragraph{As noted in Chapter~\ref{chap:Step 1 — Identifying Hazards and Risk Sources}, people tend to underestimate the likelihood of negative events affecting them personally, even when they have accurate statistical knowledge about base rates. In a communication context, this means that audiences will often discount risk assessments that relate to their own organisation or their own area of responsibility — applying an implicit assumption that \textit{"that kind of thing doesn't happen to organisations like us."}.}
\paragraph{Effective risk communication counters this by using specific, comparable examples — peer organisations, similar processes, equivalent contexts — that make it harder for the audience to maintain the comfortable fiction that the risk is someone else's problem.}

\subsection{Dread Risk}
\paragraph{Some categories of risk generate a disproportionate emotional response — not because they are statistically the most significant but because they are perceived as particularly uncontrollable, catastrophic, or threatening to fundamental values. Cybersecurity risk, reputational damage, personal liability for senior individuals, and the possibility of criminal prosecution all tend to generate this kind of heightened response.}
\paragraph{Understanding which risks carry this emotional loading for your specific audience — and calibrating your communication accordingly — is part of effective risk communication. It means neither amplifying dread risks out of proportion to their actual significance, nor treating them with the same detached analytical tone that might be appropriate for a more routine operational risk.}

%---------------------------- SECTION ----------------------------
\section{The Medium Matters — Choosing the Right Format}
\paragraph{The content of a risk communication is only part of the picture. The format and medium through which it is delivered significantly affects how it is received and how much of it is retained.}
\paragraph{\textbf{Written reports} — the most common format for formal risk communication — are appropriate for detailed analysis, for creating a documented record, and for audiences who need the ability to refer back to the information over time. They work less well for complex or nuanced messages that benefit from discussion, for audiences who are unlikely to read dense written material carefully, or for communications that need to generate immediate engagement and response.}
\paragraph{\textbf{Risk dashboards and visual summaries} — heat maps, trend charts, KRI dashboards, and other visual formats — are powerful tools for conveying the overall shape of the risk picture quickly and memorably. They work well for board-level summaries and management information, where the goal is to give a rapid, accessible overview rather than detailed analysis. They work less well for communicating nuance, context, or the reasoning behind assessments.}
\paragraph{\textbf{Presentations and discussions} — oral communication, whether in board meetings, risk committees, or management briefings, allows for real-time engagement, question and answer, and the kind of iterative dialogue that written reports cannot facilitate. For important or complex risk communications, the combination of a concise written summary with a focused oral presentation — designed to drive genuine discussion rather than simply read out the document — is often the most effective approach.}
\paragraph{\textbf{Training and awareness programmes} — for operational risk communication, structured learning and awareness activities are often more effective than formal reports or presentations. Well-designed training helps people understand risks in the context of their own work, builds genuine competence in risk identification and management, and reinforces the risk-aware culture discussed in the next chapter.}
\paragraph{\textbf{Informal communication} — the conversations that happen outside formal governance structures — in management meetings, in one-to-one discussions, in the daily interactions between compliance professionals and their business colleagues — are as important as formal risk communications, and sometimes more so. The compliance professional who is present, engaged, and approachable — who is part of the conversation rather than a distant producer of formal reports — is far more likely to be brought into risk discussions early, when there is still time to act, than one whose presence is felt only through formal channels.}

%---------------------------- SECTION ----------------------------
\newpage
\section{Chapter Summary}
In this chapter we learned about the following key points:
\begin{table}[H]
  \centering
  \renewcommand{\arraystretch}{1.5}
  \begin{tabularx}{\textwidth}{ | >{\raggedright\arraybackslash}p{0.20\textwidth} 
								| >{\raggedright\arraybackslash}p{0.40\textwidth} 
                                | >{\raggedright\arraybackslash}X | }
    \hline
    \textbf{Challenge} & \textbf{Why It Matters} & \textbf{How to Address It} \\
    \hline
    \textbf{Risk is abstract} & People struggle to engage with things that have not happened & Make risks concrete through scenarios, examples, and analogies\\
    \hline
    \textbf{Numbers without narrative} & Data without context communicates little & Provide insight, not just information — explain what the data means\\
    \hline
    \textbf{Organisational dynamics} & Informal pressures suppress bad news & Build credibility, frame constructively, and have the courage to deliver difficult messages\\
    \hline
    \textbf{Audience diversity} & Different people need different things & Tailor content, detail, and format to each specific audience\\
    \hline
    \textbf{Optimism bias} & People discount risks to their own organisation & Use peer examples and comparable cases to counter comfortable assumptions\\
    \hline
  \end{tabularx}
  \caption{Communication Challenges}
\end{table}

\begin{table}[H]
  \centering
  \renewcommand{\arraystretch}{1.5}
  \begin{tabularx}{\textwidth}{ | >{\raggedright\arraybackslash}p{0.35\textwidth} 
								| >{\raggedright\arraybackslash}p{0.35\textwidth} 
                                | >{\raggedright\arraybackslash}X | }
    \hline
    \textbf{Audience} & \textbf{Focus} & \textbf{Key Characteristics} \\
    \hline
    \textbf{Board} & Strategic overview, material risks, appetite alignment & Concise, forward-looking, honest, invites genuine engagement\\
    \hline
    \textbf{Senior management} & Integrated risk picture, cross-functional issues, action requirements & Cross-functional, action-oriented, genuinely challenging\\
    \hline
    \textbf{Operational teams} & Specific, practical, directly relevant to their work & Concrete, two-way, supported by training\\
    \hline
    \textbf{Regulators} & Accurate, transparent, consistent & Honest, proactive, demonstrating genuine engagement\\
    \hline
  \end{tabularx}
  \caption{Communication Audiences}
\end{table}

\paragraph{Key takeaways:}
\begin{itemize}
  \item{Risk communication is not the same as distributing risk information — it requires genuine effort to ensure that the right people develop genuine understanding.}
  \item{The challenges of risk communication are rooted in human psychology and organisational dynamics, not just technical complexity.}
  \item{Effective risk communication leads with significance, makes the abstract concrete, provides insight rather than information, and makes the call to action explicit.}
  \item{Different audiences require fundamentally different approaches — tailoring communication to the specific audience is not optional, it is essential.}
  \item{Honesty is non-negotiable — risk communications that soften or suppress difficult messages fail at the most basic level of their purpose.}
  \item{The medium matters — the right format for the message and the audience significantly affects how it is received.}
  \item{Informal communication is as important as formal reporting — the compliance professional's presence and accessibility within the organisation is a genuine risk communication asset.}
\end{itemize}

\barquote{In Chapter~\ref{chap:Building a Risk-Aware Culture}, we build on the communication themes of this chapter to explore the broader question of risk culture — what it is, why it matters, how it is shaped, and what compliance professionals can do to support the development of a genuinely risk-aware organisation.}